\documentclass[11pt,a4paper]{article}
\usepackage[T1]{fontenc}
\usepackage{lmodern,amsmath,amssymb}
\usepackage[margin=23mm,top=19.5mm,bottom=22mm]{geometry}
\usepackage{microtype}
% Keep a rebuild independent of the host clock and absolute source path.
\pdfinfoomitdate=1
\pdftrailerid{}
\pdfsuppressptexinfo=15
\usepackage[hidelinks,pdflang={en},
  pdfauthor={Astra 6.2},
  pdftitle={Undecidability of freeness for finitely generated subalgebras},
  pdfsubject={Formalized in Lean by Opus 5.5. Supervised by Christophe Cordero.
    Discovered on 26 November 2026.}]{hyperref}
\newcommand{\Ld}{\mathsf{L}}
\newcommand{\Rd}{\mathsf{R}}
\newcommand{\blank}{\square}
\newcommand{\head}[1]{\par\noindent\textbf{#1}\enspace}
\setlength{\parskip}{.2em}
\begin{document}
\begin{center}
  {\LARGE\bfseries Undecidability of freeness for\\[2pt]finitely generated subalgebras\par}
\end{center}
\noindent\begin{tabular*}{\linewidth}{@{}l@{\extracolsep{\fill}}lll@{}}
{\small\sffamily Author} & {\small\sffamily Lean formalization} & {\small\sffamily Supervisor} &
  {\small\sffamily Date of discovery} \\
Astra 6.2 & Opus 5.5 & Christophe Cordero & 26 November 2026
\end{tabular*}
\par\smallskip\hrule\medskip

\head{Theorem.}Fix a field $K$. Even for a finite set $Y\subseteq K\langle a,b\rangle$ specified
by words and differences of two words, no algorithm decides whether the unital $K$-subalgebra
generated by $Y$ is isomorphic to $K\langle B\rangle$ for some alphabet $B$. Thus the problem for
arbitrary finite sets of polynomials is undecidable over every computable field. The conclusion
also holds for ring isomorphisms.

\head{Proof.}The words on an alphabet, including the empty word $1$, form a basis of its free
algebra, with multiplication by concatenation. Encode any finite auxiliary alphabet
$c_1,\dots,c_n$ by $c_j\mapsto a^jb$. Each $b$ ends one block, so distinct words have distinct
images. Linear extension embeds the free algebra in $K\langle a,b\rangle$ and preserves
differences of words.

Ring and $K$-algebra freeness agree here. Choose a top-degree word in each of two nonzero
polynomials $f,g$. Their concatenation comes only from this pair, since the split after $\deg f$
letters is forced. Thus degrees add, so the units of every unital $K$-subalgebra are exactly
$K^\times$. A ring isomorphism from $K\langle B\rangle$ onto such a subalgebra restricts to an
automorphism of $K=K^\times\cup\{0\}$. Precomposing with its inverse on coefficients, fixing the
letters of $B$, gives a $K$-algebra isomorphism.

\head{A graph criterion.}Let $D,S$ be disjoint finite alphabets and $\ell_i\to r_i$ a finite list
of rules with both sides in $D^*SD^*$. Take $D^*SD^*$ as the vertex set of a directed multigraph.
For each rule and each $u,v\in D^*$, put an edge $(u,i,v)$ from $u\ell_iv$ to $ur_iv$. Let
$A\subseteq K\langle D\cup S\rangle$ be generated by $D$ and the differences $\ell_i-r_i$. If no
undirected cycle uses an edge $(1,i,1)$, then
\[
  A \text{ is a free $K$-algebra} \quad\iff\quad
  \text{the underlying undirected multigraph is a forest.}
\]
Cycles here are finite. A loop is a cycle of length one, and two parallel edges give a cycle of
length two.

Let $Z=\{z_i\}$ contain one new letter per rule. The surjection
$\pi:K\langle D\cup Z\rangle\to A$ fixes $D$ and sends $z_i$ to $\ell_i-r_i$. It preserves degree
when $D$ has degree zero and $Z,S$ have degree one. In degree one it is the incidence map
\[
  uz_iv \longmapsto u\ell_iv-ur_iv.
\]
A signed sum around a cycle is a nonzero kernel element over every field. An edge on no cycle
has coefficient zero in every kernel element. Deleting it disconnects its endpoints, and summing
the vertex coefficients on one side gives, up to sign, its coefficient. All sums have finite
support.

In a forest the degree-one map is therefore injective. Call the letters of $S\cup Z$
distinguished. In degree $k\ge1$, replace the $k$ rule letters from left to right. At position
$j$, fix everything through the preceding distinguished letter and from the following one onward,
using an empty prefix or suffix when there is no such letter. In the middle the substitution is
exactly the degree-one map. Distinct fixed contexts have disjoint output supports because the
$k$ distinguished letters remain visible. Each step is injective, and $\pi$ fixes degree zero.
Hence $A$ is free.

Let $I$ be the zero-constant-term ideal of $K\langle D\cup Z\rangle$. For any relation, its
degree-zero part vanishes since $\pi$ fixes $K\langle D\rangle$. Its degree-one part has zero
coefficient at each bare $z_i$ by the hypothesis on $(1,i,1)$. Thus every relation has zero
constant and linear terms in ordinary word length, giving $\ker\pi\subseteq I^2$.

Put $J=\pi(I)$, the kernel of the constant-term map on $A$. Since $J^2=\pi(I^2)$, the induced
map $I/I^2\to J/J^2$ is an isomorphism. If $A$ is free, subtract the constant term from each free
generator. This invertible translation gives free coordinates in which $J/J^2$ is the space of
linear forms. The linear parts of the images of $D\cup Z$ are therefore independent. The linear
map sending these forms back to their corresponding letters and a complementary basis to zero
extends multiplicatively to a left inverse. Hence the lowest nonzero ordinary-degree component of
a polynomial survives evaluation in its own degree. Freeness forces $\ker\pi=0$, which a cycle
precludes. This proves the criterion.

\head{Encoding a computation.}Take a deterministic Turing machine with a two-way infinite tape,
alphabet $\Gamma$, blank $\blank$, initial state $q_0$, and a transition for every state-symbol
pair outside its halting state $h$. Fix an input word $w$, using $w=\blank$ for empty input. Set
\[
  D=\Gamma\;\sqcup\;\{\bar c:c\in\Gamma\}\;\sqcup\;\{\Ld,\Rd\}.
\]
The letters $\Ld,\Rd$ delimit a finite tape segment. The separate alphabet $S$ contains
$H_q,P_q,U_q,F_q,B_q$ for each state $q$, and two cleanup letters $C,E$. The initial word is
$\Ld H_{q_0}w\Rd$, with the scanned cell immediately after the control letter.

For each instruction in state $q$ reading $c$, writing $d$ and entering state $r$, include the
appropriate rule
\[
  \text{right move}\quad H_qc\to dP_r,\qquad\qquad \text{left move}\quad H_qc\to U_rd.
\]
Include the following rules for every state $r$ and plain tape symbol $c\in\Gamma$, taking each
distinct rule once.
\[
  \begin{array}{lll}
    cU_r\to P_rc, & \Ld U_r\to \Ld P_r\blank, & P_r\Rd\to P_r\blank\Rd,\\[3pt]
    P_rc\to \bar cF_r, & F_rc\to cF_r, & F_r\Rd\to B_r\blank\Rd,\\[3pt]
    cB_r\to B_rc, & \bar cB_r\to H_rc. &
  \end{array}
\]
Here $U_r$ completes a left move, and $P_r$ marks the new head cell. The end rules supply a blank
when that cell is outside the represented segment. The $F_r$ sweep adds a blank at the right
delimiter and $B_r$ returns to the marked cell and restores it. Starting from the initial word,
each such round therefore simulates one instruction exactly.

For cleanup use the following rules, with $x$ ranging over plain and marked tape symbols.
\[
  \begin{array}{lll}
    H_h\to C, & Cx\to xC, & C\Rd\to E\Rd,\\[3pt]
    xE\to E, & \Ld E\to \Ld H_{q_0}w. &
  \end{array}
\]
After a halt, $C$ reaches the right delimiter and $E$ erases the tape back to the left delimiter.
The last rule then restores the initial word.

\head{Cycles and termination.}Even on an arbitrary word with one control letter, at most one rule
applies. Its control letter fixes the direction of inspection, and each inspected symbol selects
at most one rule. In any directed multigraph of outdegree at most one, every undirected cycle is
directed. Indeed its edges use one outgoing edge per vertex on average, hence exactly one at
every vertex.

Every directed cycle in this rule system uses the reset $\Ld E\to \Ld H_{q_0}w$. To see this, a
path in the cleanup phase can only move $C$ to the right and then erase with $E$, until it stops
or resets. Outside cleanup, length never decreases. The preparatory moves take finitely many
steps and each individual sweep moves in one direction along a finite word. A cycle through the
simulation phases would therefore use $F_r\Rd\to B_r\blank\Rd$, which strictly increases the
length.

On a cycle, erasure before a reset began with $C\Rd\to E\Rd$. The symbol immediately following
$E$ stays $\Rd$ throughout erasure. The reset therefore produces $u\Ld H_{q_0}w\Rd v$ for some
exterior words $u,v$. All further moves act between these same delimiters, so the next cleanup
requires the given machine to halt on $w$. Conversely, halting returns the initial word to
itself in a positive number of steps. Thus a cycle exists exactly when the machine halts.

Every cycle contains both delimiters. All rules preserve their two counts separately, whereas
each rule with empty exterior contexts contains at most one delimiter. Its edge therefore lies in
a component containing no cycle. The graph criterion now gives
\[
  A\text{ is free}\quad\iff\quad\text{the machine does not halt on $w$.}
\]

The finite rule list and its binary encoding are effective. A freeness decider would therefore
decide halting. If a halting decider existed, a machine could copy its input $e$, run the decider
on $(e,e)$, loop on a positive answer, and halt on a negative one. Taking $e$ to be that machine's
own binary description contradicts either answer.\hfill$\square$
\end{document}
